% ============ 第4章 §4.1 Cotlar 引理 (原书 pp.75-76) ============

\section{Cotlar 引理}\label{sec:4.1}

这一节在抽象的框架下展开讨论. 主要目的是处理奇异积分算子, 特别是 Hilbert 变换, 但在 Hilbert 空间的框架下作抽象的讨论.

这里证明后面要使用的正交性引理. 对 Hilbert 空间 $H$, $P(\mathcal{H})$ 表示全体正交射影组成的集合. 即, 有界线性算子 $P: H \to H$ 属于 $P(\mathcal{H})$, 是指 $P = P^{*}$ 与 $P^{2} = P$ 成立. 这里, 对不同的自然数 $j, k$, 对满足正交条件 $P_{j}P_{k} = 0$ 的正交射影列 $\{P_{j}\}_{j=1}^{\infty} \subset P(\mathcal{H})$,
\begin{equation}\tag{4.1}
\|P_{1} + P_{2} + \cdots + P_{N}\|_{\mathcal{H} \to \mathcal{H}} \leq 1
\label{eq:4.1}
\end{equation}
成立. 由于极限在适当的意义下存在,
\[
\|P_{1} + P_{2} + \cdots\|_{\mathcal{H} \to \mathcal{H}} \leq 1
\]
也成立. 但是, 像 $P_{1} = P_{2} = \cdots = A$ 的情形那样, 对不是正交射影的一般线性算子, 式~\eqref{eq:4.1} 不成立.

从应用上看, 这个正交条件经常出现, 但这个关系式只要稍加摄动就不再成立. 这里, 寻找正交条件的替代.

\begin{theorem}[Cotlar 引理]\label{thm:4.1}
设 $J \subset \mathbb{Z}$ 为有限集合. 对有界线性算子的族 $\{T_{j}\}_{j \in J} \subset B(\mathcal{H})$, 存在适当的正数列 $\{\gamma_{j}\}_{j \in \mathbb{Z}}$, 使得
\begin{equation}\tag{4.2}
A \equiv \max\left(\sup_{j' \in J}\sum_{j \in J}\gamma_{j-j'}, \sup_{j' \in J}\sum_{j \in J}\gamma_{j'-j}\right) < \infty
\label{eq:4.2}
\end{equation}
以及
\begin{equation}\tag{4.3}
\|T_{j'}^{*}T_{j}\|_{\mathcal{H} \to \mathcal{H}} + \|T_{j'}T_{j}^{*}\|_{\mathcal{H} \to \mathcal{H}} \leq \gamma_{j'-j}^{2} \quad (j, j' \in J)
\label{eq:4.3}
\end{equation}
成立. 此时, 算子 $T_{J} \equiv \sum_{j \in J}T_{j}$ 满足 $\|T_{J}\|_{\mathcal{H} \to \mathcal{H}} \leq A$.
\end{theorem}

\begin{proof}
设自然数 $m$ 满足 $\log_{2}m \in \mathbb{N}$. 认真地写下来的话,
\begin{equation}\tag{4.4}
(T_{J}^{*}T_{J})^{m} = \sum_{j_{1}, j_{2}, \ldots, j_{2m} \in J}T_{j_{1}}^{*}T_{j_{2}}T_{j_{3}}^{*} \cdots T_{j_{2m}}
\label{eq:4.4}
\end{equation}
成立.

评价各项. 首先, 将括号如下配对.
\[
T_{j_{1}}^{*}T_{j_{2}}T_{j_{3}}^{*} \cdots T_{j_{2m}} = \prod_{l=1}^{m}\left(T_{j_{2l-1}}^{*}T_{j_{2l}}\right).
\]
由不等式~\eqref{eq:4.3},
\begin{equation}\tag{4.5}
\|T_{j_{1}}^{*}T_{j_{2}}T_{j_{3}}^{*} \cdots T_{j_{2m}}\|_{\mathcal{H} \to \mathcal{H}} \leq \prod_{l=1}^{m}\gamma_{j_{2l-1}-j_{2l}}^{2}
\label{eq:4.5}
\end{equation}
成立.

接下来, 将括号如下配对.
\[
T_{j_{1}}^{*}T_{j_{2}}T_{j_{3}}^{*} \cdots T_{j_{2m}} = T_{j_{1}}^{*}\left(\prod_{l=1}^{m-1}\left(T_{j_{2l}}T_{j_{2l+1}}^{*}\right)\right)T_{j_{2m}}.
\]
这里, 关于两端, 由式~\eqref{eq:4.3},
\[
\|T_{j_{1}}\|_{\mathcal{H} \to \mathcal{H}} = \sqrt{\|T_{j_{1}}^{*}T_{j_{1}}\|_{\mathcal{H} \to \mathcal{H}}} \leq \gamma_{0} \leq A
\]
成立. 同样地, 由 $\|T_{j_{2m}}\|_{\mathcal{H} \to \mathcal{H}} \leq A$,
\begin{equation}\tag{4.6}
\|T_{j_{1}}^{*}T_{j_{2}}T_{j_{3}}^{*} \cdots T_{j_{2m}}\|_{\mathcal{H} \to \mathcal{H}} \leq A^{2}\prod_{l=1}^{m-1}\gamma_{j_{2l}-j_{2l+1}}^{2}
\label{eq:4.6}
\end{equation}
成立. 考虑式~\eqref{eq:4.5} 与式~\eqref{eq:4.6} 的几何平均, 代入式~\eqref{eq:4.4},
\[
\|(T_{J}^{*}T_{J})^{m}\|_{\mathcal{H} \to \mathcal{H}} \leq A\sum_{j_{1}, j_{2}, \ldots, j_{2m} \in J}\prod_{l=1}^{2m-1}\gamma_{j_{l}-j_{l+1}}
\]
成立.

这里, 先对 $j_{1}$ 求和, 利用 $\sum_{j_{1} \in J}\gamma_{j_{1}-j_{2}} \leq A$ 来控制. 接下来对 $j_{2}$ 求和, 利用 $\sum_{j_{2} \in J}\gamma_{j_{2}-j_{3}} \leq A$ 来控制. 如此反复, 只特别对待 $j_{2m}$, 得 $\|T_{J}\|_{\mathcal{H} \to \mathcal{H}}^{2m} = \|(T_{J}^{*}T_{J})^{m}\|_{\mathcal{H} \to \mathcal{H}} \leq A^{2m}\sharp J$. 这里, 由于 $J$ 是有限的, 得 $\|T_{J}\|_{\mathcal{H} \to \mathcal{H}} \leq A(\sharp J)^{\frac{1}{2m}}$. 最后令 $m \to \infty$, 即得结论.
\end{proof}

至此 $J$ 一直取有限集合, 但对 $J = \mathbb{Z}$ 定理也能适用. 此时 $T_{J} = T_{\mathbb{Z}}$ 是无穷和, 但可以以对一切 $x \in H$,
\[
T_{\mathbb{Z}}x = \sum_{j \in \mathbb{Z}}T_{j}x
\]
收敛的意义来理解无穷和. 实际上, 若存在 $j$ 使得 $x \in \mathrm{Im}(T_{j}^{*})$, 则定义无穷和的极限的存在性由式~\eqref{eq:4.2} 与式~\eqref{eq:4.3} 是显然的. 由于在有限和的情形已证明算子范数具有不依赖于 $J$ 的估计, 移到线性闭包,
\[
x \in \overline{\operatorname{Span}}\left(\mathrm{Im}(T_{j}^{*}), j \in \mathbb{Z}\right)
\]
的情形, 极限也同样存在. 最后, 考虑
\[
x \in \bigcap_{j \in \mathbb{Z}}\ker(T_{j}) = \overline{\operatorname{Span}}\left(\mathrm{Im}(T_{j}^{*}), j \in \mathbb{Z}\right)^{\perp}
\]
的情形. 这里 Span 是在代数的意义下使用的表示线性空间的生成的记号. 即, 若 $T_{j}x = 0$ 对一切 $j \in \mathbb{Z}$ 成立, 则由对一切有限集合 $J \subset \mathbb{Z}$ 有 $T_{J}x = 0$, 极限是存在的.

\begin{exercise}\label{exer:4.1}
用 $\|\cdot\|$ 也表示复 Hilbert 空间 $H$ 中全体有界线性算子组成的 Banach 空间 $B(H)$ 的范数. 此时, 证明 $(A^{*}A)^{2^{m}}$ 是自伴的, 且 $\|(A^{*}A)^{2^{m}}\| = \|A\|^{2^{m+1}}$ 成立.
\end{exercise}
